\documentclass[11pt]{article}
\usepackage{amsmath}
\usepackage{amssymb}
\usepackage{geometry}
\geometry{margin=1in}

\title{Quintic Doc}
\author{}
\date{}

\begin{document}
\maketitle
\vspace{-3em}


Summary: At each control cycle/step, the desired pose at a look-ahead time is computed to obtain error twist which is then mapped to joint velocities using differential inverse kinematics.

\section*{Quintic Polynomial for Position}

Let $p(t)\in\mathbb{R}$ denote end-effector position.  
A quintic trajectory has the form
\[
p(t)=a_0 + a_1 t + a_2 t^2 + a_3 t^3 + a_4 t^4 + a_5 t^5.
\]

(Quintic polynomial was chosen based on a prior team meeting where we discussed it will be good to have smooth acceleration profiles (S-shaped).)

To guarantee continuous position, velocity, and acceleration, the boundary conditions
\[
p(0)=p_0,\quad p(T)=p_T,\quad 
\dot p(0)=v_0,\ \dot p(T)=v_T,\quad
\ddot p(0)=a_0,\ \ddot p(T)=a_T
\]
yield the linear system $A\mathbf{a}=\mathbf{b}$, from which the coefficients $\mathbf{a}=[a_0,\dots,a_5]^\top$ are solved.  
The translational trajectory in $\mathbb{R}^3$ is obtained by fitting three independent quintics to the $x,y,z$ coordinates of the horizon points. (There should be a way to vectorize this instead of doing one axis at a time like the current code does.)

The instantaneous desired position and linear velocity follow directly:
\[
p(t) = \sum_{i=0}^5 a_i t^i,\qquad
\dot p(t) = \sum_{i=1}^5 i\,a_i t^{i-1}.
\]

\section*{SLERP for Orientation}

At the same timestep $at\_time\_t$ that position is computed, orientation is obtained via spherical linear interpolation (SLERP). This is pretty standard and there's not much to add in here. I used scipy's SLERP function.

\section*{Pose Error}

The desired pose at timestep $at\_time\_t$ is constructed from the position and orientation computed above. Error twist is obtained by comparing the desired pose to the current end-effector pose.

Staying consistent with $cartesianDualArmPathFollowingControlLoop()$, I included the multiplication by 2 for the orientation. I actually wanted to bring up why this was needed, my best guess is empirical tests revealed this scaling to be needed.

\section*{Differential Inverse Kinematics}

This is exactly the same as in $cartesianDualArmPathFollowingControlLoop()$. I also ensured the same return type so that my implementation can be a drop-in replacement.

\end{document}
